\subsection{Arithmetic approximation in finite transcendence degree}

\begin{lemma}[Proved; the affine bound over an arithmetic curve]
\label{A:arithmetic-affine-relative}
Let $F$ be a number field, let $p$ be a prime, and let
$S=\Spec\mathcal O_F[1/p]$. Suppose that $p$ is odd or that
$F$ is totally imaginary. Let $f:T\to S$ be an affine, flat
morphism of finite type whose generic fibre has dimension at
most $e$. Then
\[
 \cd_p(T)\leq e+2.
\]
The scheme $T$ need not be integral or normal.
\end{lemma}
\begin{proof}
The base has $p$-cohomological dimension at most two, by
Artin--Verdier duality in the form used in
Proposition~\ref{A:arithmetic-test}. For odd $p$ the real-place
terms have no positive $p$-cohomology; thus the same upper bound
holds without the totally imaginary hypothesis. Every closed
point of $S$ has a finite residue field and hence
$p$-cohomological dimension one.

Give $S$ its dimension function $\delta_S$ with value one at
the generic point and zero at every closed point. Flatness
excludes vertical irreducible components of $T$. The dimension
formula for finite-type schemes over the excellent regular
one-dimensional scheme $S$ shows that all fibres have dimension
at most $e$. Therefore the induced function
\[
 (f^*\delta_S)(t)
 =\delta_S(f(t))+\operatorname{trdeg}_{\kappa(f(t))}\kappa(t)
\]
has supremum at most $e+1$. This argument applies to every
irreducible component and makes no integrality assumption on $T$.

Let $\mathcal F$ be a constructible sheaf of
$\mathbf F_p$-vector spaces on $T$. Gabber's affine theorem
\citep[Expos\'e XV, Theorem 1.1.2 and its reformulation 1.2.1,
p.~219]{A_ILO2014} gives
\[
 \delta_S\bigl(R^b f_*\mathcal F\bigr)\leq e+1-b.
\]
Thus $R^b f_*\mathcal F=0$ for $b>e+1$, and
$\mathcal Q=R^{e+1}f_*\mathcal F$ has zero geometric generic
stalk. The latter conclusion can also be read from the ordinary
affine theorem on the geometric generic fibre.

We check the cohomology of this top direct image using the
generic-zero sheaf argument in the proof of
Theorem~\ref{A:curve-normalization}. Write
$\mathcal Q=\varinjlim_i\mathcal Q_i$ with each
$\mathcal Q_i$ constructible. For each $i$, finiteness of its
generic stalk and vanishing of the generic colimit give a later
$j$ for which $\mathcal Q_i\to\mathcal Q_j$ is zero at the
generic point. The image of this transition is constructible
and supported on finitely many closed points. Exact closed
pushforward and the dimension-one bound for finite fields give
vanishing above degree one for this image. The transition on
cohomology factors through that image, so every class vanishes
at a later stage. Cohomological continuity gives
$H^a(S,\mathcal Q)=0$ for $a>1$, whether or not constructibility
of $\mathcal Q$ is invoked.

In the Leray spectral sequence
\[
 H^a\bigl(S,R^b f_*\mathcal F\bigr)
 \Longrightarrow H^{a+b}(T,\mathcal F),
\]
the rows $b\leq e$ vanish for $a>2$, the row $b=e+1$ vanishes
for $a>1$, and higher rows vanish identically. Hence the total
cohomology vanishes in degrees greater than $e+2$.
Coefficient d\'evissage and constructible approximation on the
noetherian scheme $T$ prove the bound for all $p$-primary torsion
sheaves.
\end{proof}

\begin{proposition}[Proved; arithmetic approximation beyond finite type]
\label{A:arithmetic-transcendence}
Let $R$ be a characteristic-zero domain, let $p$ be invertible
in $R$, and suppose
$d=\operatorname{trdeg}_{\Q}\operatorname{Frac}R<\infty$.
No noetherian hypothesis is imposed. If $p$ is odd, then
\[
 \cd_p(\Spec R)\leq d+2.
\]
If $p=2$ and $K=\operatorname{Frac}R$ is not orderable, there
is a nonzero $a\in R$ such that
\[
 \cd_2\bigl(\Spec R[1/a]\bigr)\leq d+2.
\]
The same bounds apply to every affine open of the indicated scheme.
\end{proposition}
\begin{proof}
For odd $p$, write $R$ as the filtered union of its finite-type
$\Z[1/p]$-subalgebras $R_i$. These are domains and are flat over
$\Z[1/p]$, since they are torsion free as modules over a
Dedekind domain. Their generic fibres have dimensions
$e_i=\operatorname{trdeg}_{\Q}\operatorname{Frac}(R_i)\leq d$.
Lemma~\ref{A:arithmetic-affine-relative} gives
$\cd_p(R_i)\leq d+2$ uniformly. Descent of constructible
coefficients, continuity of cohomology, and the passage to all
torsion sheaves in Lemma~\ref{A:A2-approximation} give the same
bound on $R$.

Now let $p=2$. Write $-1=\sum_{j=1}^m x_j^2$ in $K$ and
choose a common nonzero denominator $a$ so that each
$x_j\in R[1/a]$. Approximate $R[1/a]$ by its finite-type
$\Z[1/2]$-subalgebras $A_i$ containing all $x_j$.
Every residue field of $A_i$ is nonorderable, because the displayed
identity survives in it. The fraction field of $A_i$ is a
finitely generated nonorderable characteristic-zero field. It has
finite $2$-cohomological dimension, by the field theorem
\citep[Expos\'e X, Theorem 2.1]{A_SGA4}. Consequently
Corollary~\ref{A:generic-plus-dimension} gives
$\cd_2(A_i)<\infty$.

Use the quadratic finite \emph{\'etale cover}
\[
 A_i'=A_i[T]/(T^2+1),\qquad
 \Spec A_i'\longrightarrow\Spec A_i.
\]
We do not assume that $A_i'$ is a domain. As a scheme over
$\Spec\Z[i,1/2]$, it is the base change of the flat
$\Spec A_i\to\Spec\Z[1/2]$. It is therefore affine, flat
and of finite type, and its generic fibre has dimension at most
$d$. Lemma~\ref{A:arithmetic-affine-relative}, with the totally
imaginary number field $\Q(i)$, gives $\cd_2(A_i')\leq d+2$.
The cover is a torsor under the constant group of order two, so
Lemma~\ref{A:quadratic-descent}, using the finite bound already
proved for $A_i$, gives $\cd_2(A_i)\leq d+2$.
The same continuity argument passes this uniform estimate to
$R[1/a]$.

Finally, the coordinate ring of a nonempty affine open has the
same fraction field. For odd $p$ the first argument applies to it
directly. In the second case it contains the elements $x_j$, so
the approximation by subalgebras containing the one fixed
sum-of-squares identity applies without a further localization.
This proves the last assertion.
\end{proof}

\begin{corollary}[Proved; arbitrary noetherian models of maximal arithmetic dimension]
\label{A:A2-arithmetic-maximal}
Let $R$ be a noetherian characteristic-zero domain, let $p$ be
invertible in $R$, and put $K=\operatorname{Frac}R$. If
\[
 \cd_p(K)=\operatorname{trdeg}_{\Q}K+2<\infty,
\]
then A.2 holds for $\Spec R$. In particular it holds for
\emph{every} noetherian model whose fraction field is finitely
generated over $\Q$ and whose $p$-cohomological dimension is
finite. Excellence, normality and finite type of the model
are unnecessary.
\end{corollary}
\begin{proof}
For $p=2$ finite cohomological dimension excludes an ordering of
$K$. Proposition~\ref{A:arithmetic-transcendence} therefore
supplies an open $V$ whose every affine open has cohomological
dimension at most $n=\cd_p(K)$; for odd $p$ we may take
$V=\Spec R$. Given a nonempty open $W$, apply
Proposition~\ref{A:spread} inside $W\cap V$. Its lower witness
survives on every smaller nonempty affine open, giving equality
there. This is A.2 with its full quantifier over smaller opens.

For a field finitely generated over $\Q$, the field theorem
\citep[Expos\'e X, Theorem 2.1]{A_SGA4} gives the displayed
equality for odd $p$. For $p=2$, apply it after adjoining $i$;
finite extension preserves finite $2$-cohomological dimension
for the nonorderable field $K$. Thus the last assertion has exactly
the stated hypothesis.
\end{proof}

\begin{remark}[The remaining arithmetic range]
The bound is obtained from the opened affine and arithmetic
duality theorems; we make no priority claim. We searched the
phrases ``affine cohomological dimension arithmetic schemes'',
``cohomological dimension finite type over integers'', and
``dimension cohomologique affine arithm\'etique'', and checked
SGA~4, Expos\'es X and XIV, and the precise affine statement in
\citet[Expos\'e XV, Theorem 1.1.2]{A_ILO2014}.
This reduction leaves, among finite-transcendence-degree
characteristic-zero fields, those for which the cohomological
dimension drops below $\operatorname{trdeg}_{\Q}K+2$.
The cyclotomic arithmetic instance remains in that range and
is not settled by this approximation argument.
\end{remark}
